\documentclass[a4paper]{article}     
\usepackage{amsfonts}
\usepackage{amsmath}
\usepackage{amssymb}
\usepackage{graphicx}

\begin{document}

\noindent \large Auteur: Mohamed Shafik \\
\noindent \large Studierichting: Informatica\\
\noindent \large Oefeningenreeks: 6B nr 14\\

\medskip

\normalsize

De som van de eerste $n$ natuurlijke getallen is 780. Bepaal $n$.\\

Rekenkundige rij met $x_1 = 1$, $x_n = n$, $s_n = \dfrac{x_1 + x_n}{2} \cdot n$\\

\begin{eqnarray*}
s_n = \dfrac{x_1 + x_n}{2} \cdot n = \dfrac{1 + n}{2} \cdot n = \dfrac{n(n+1)}{2} &=& 780\\
n(n+1) &=& 1560\\
n^2 + n - 1560 &=& 0\\
\end{eqnarray*}

Berekening discriminant\\

\begin{eqnarray*}
D &=& 1^2 + 4 \cdot 1560 = 1 + 6240 = 6241\\
\sqrt{D} &=& 79\\
\end{eqnarray*}

Berekening $n$ met (abc-formule)\\

\begin{eqnarray*}
n &=& \dfrac{-1 + 79}{2} = \dfrac{78}{2} = 39\\
\end{eqnarray*}

($n = \dfrac{-1 - 79}{2} = -40$ voldoet niet want $n \in \mathbb{N}$)\\

Verificatie: $s_{39} = \dfrac{1 + 39}{2} \cdot 39 = \dfrac{40}{2} \cdot 39 = 780$\\

\begin{thebibliography}{999}
\end{thebibliography}
\end{document}
